% संपूर्ण मराठी ग्रंथातील प्रकरण 17: उपपत्ती आणि त्यांची प्रतिमाने
% हा संपादनयोग्य स्रोतखंड आहे. संकलनासाठी संपूर्ण स्रोतसंग्रहातील
% अक्षररूपे, पूर्वभाग, चिन्हव्याख्या आणि आकृती-स्रोत आवश्यक आहेत.
% मूळ: Open Logic Project; मजकूर आणि रूपांतर CC BY 4.0.
% यंत्रानुवाद, दुरुस्ती आणि तपासणी: OpenAI Codex —
% GPT-5.6 Sol आणि GPT-6 Sol, दोन्ही Ultra effort.
% दुरुस्ती-पुनरावलोकन: GPT-6.1 Sol, Ultra effort.
\chapter{उपपत्ती आणि त्यांची प्रतिमाने}\label{fol:mat::chap}








\section{परिचय}\label{fol:mat:int:sec}

\begin{explain}
स्वयंसिद्धकीय पद्धतीचा विकास हा विज्ञानाच्या इतिहासातील एक महत्त्वपूर्ण
पराक्रम आहे आणि गणिताच्या इतिहासात त्याला विशेष महत्त्व आहे. एखाद्या
क्षेत्राचा स्वयंसिद्धकीय विकास करताना अनेक प्रश्न स्पष्ट करावे लागतात: हे
क्षेत्र कशाविषयी आहे? सर्वांत मूलभूत संकल्पना कोणत्या? त्यांचे परस्पर संबंध
कसे आहेत? क्षेत्रातील सर्व संकल्पना या मूलभूत संकल्पनांच्या साहाय्याने
परिभाषित करता येतात का? या संकल्पना कोणते नियम पाळतात आणि त्यांनी कोणते
नियम पाळलेच पाहिजेत?

स्वयंसिद्धकीय पद्धत आणि तर्कशास्त्र जणू एकमेकांसाठीच निर्माण झाले आहेत.
आकारिक तर्कशास्त्र स्वयंसिद्धकीय उपपत्ती मांडण्यासाठी, उपपत्तीच्या
स्वयंसिद्धकांपासून प्रमेये नेमक्या निर्दिष्ट रीतीने सिद्ध करण्यासाठी आणि
स्वयंसिद्धकांची पूर्ति करणाऱ्या सर्व प्रणालींच्या गुणधर्मांचा सुव्यवस्थित
अभ्यास करण्यासाठी साधने पुरवते.
\end{explain}

\begin{defn}
वाक्यांचा संच~$\Gamma$ \emph{बंद} असतो तेव्हा आणि केवळ तेव्हाच,
जेव्हा $\Gamma \Entails \varphi$ असल्यास $\varphi \in \Gamma$ असते.
वाक्यांच्या संच~$\Gamma$ चे \emph{संवरण} म्हणजे
$\Setabs{\varphi}{\Gamma \Entails \varphi}$ होय.

जर~$\Gamma$ हे~$\Delta$ चे संवरण असेल, तर~$\Gamma$ हा वाक्यांच्या
संच~$\Delta$ ने \emph{स्वयंसिद्धकांद्वारे निरूपित} केला आहे असे म्हणतो.
\end{defn}

\begin{explain}
एखादी स्वयंसिद्धकीय उपपत्ती म्हणजे तिच्या स्वयंसिद्धकांच्या संच~$\Delta$
ने निरूपित केलेल्या वाक्यांचा संच, असे आपण मानू शकतो. दुसऱ्या शब्दांत,
आपल्याला अभ्यासायच्या स्वयंसिद्धकीय रीतीने विकसित केलेल्या शास्त्रातील
मूलभूत संकल्पना दर्शवणारी अतार्किक चिन्हे असलेली प्रथम-क्रम भाषा आणि त्या
शास्त्राचे मूलभूत नियम व्यक्त करणारा वाक्यांचा संच दिलेला असेल,
तर स्वयंसिद्धकांपासून तार्किकरीत्या निष्पन्न होणारी या भाषेतील सर्व
वाक्ये उपपत्तीचे प्रतिनिधित्व करतात असे आपण मानू शकतो. याची व्याप्ती,
एकच आदिम संकल्पना आणि सोपी स्वयंसिद्धके असलेल्या अंशतः क्रमांच्या
उपपत्तीसारख्या साध्या उदाहरणांपासून न्यूटनीय यामिकीसारख्या गुंतागुंतीच्या
उपपत्तींपर्यंत आहे.

स्वयंसिद्धकीय पद्धतीकडे पाहण्याचा हा आकारिक दृष्टिकोन इतका महत्त्वाचा
ठरण्यामागील मुख्य तार्किक तथ्ये पुढीलप्रमाणे आहेत. $\Gamma$ ही एखाद्या
उपपत्तीची स्वयंसिद्धक-प्रणाली, म्हणजे वाक्यांचा संच, आहे असे समजा.
\begin{enumerate}
\item एखादी स्वयंसिद्धक-प्रणाली संरचनेचा अभिप्रेत वर्ग नेमका कधी
  पकडते, हे आपण काटेकोरपणे सांगू शकतो. म्हणजे, आपल्याला संरचनेच्या
  एखाद्या विशिष्ट वर्गात रस असेल, तर स्वयंसिद्धक-प्रणाली~$\Gamma$ तो वर्ग
  तेव्हा आणि केवळ तेव्हाच यशस्वीपणे पकडते, जेव्हा त्या संरचना म्हणजे
  नेमकी ती~$\Struct M$ असतात ज्यांच्यासाठी $\Sat{M}{\Gamma}$.
\item या बाबतीत आपल्याला अपयश येऊ शकते, कारण काही~$\Struct M$ साठी
  $\Sat{M}{\Gamma}$ असते, पण~$\Struct M$ ही आपण अभिप्रेत धरलेल्या
  संरचना पैकी नसते. त्यामुळे~$\Struct M$ मध्ये सत्य नसलेली
  स्वयंसिद्धके आपल्याला जोडावी लागू शकतात.
\item निदान~$\Gamma$ ही सर्व अभिप्रेत संरचनांमध्ये सत्य आहे एवढ्या
  बाबतीत आपण यशस्वी झालो, तर $\Gamma \Entails \varphi$ असताना वाक्य~$\varphi$ हे
  सर्व अभिप्रेत संरचनांमध्ये सत्य असते. म्हणून, वाक्ये सर्व
  अभिप्रेत संरचनांमध्ये सत्य आहेत हे दाखवण्यासाठी, ती
  स्वयंसिद्धकांपासून तार्किकरीत्या निष्पन्न होतात एवढे दाखवून आपण तार्किक
  साधने (उदाहरणार्थ, निष्पत्ती-पद्धती) वापरू शकतो.
\item कधीकधी आपल्या मनात अभिप्रेत संरचना नसतात; त्याऐवजी आपण
  स्वयंसिद्धकांपासूनच सुरुवात करतो: काही विशिष्ट नियमांची पूर्ति करावी अशा
  काही आदिम संकल्पना आपण घेतो आणि ते नियम स्वयंसिद्धक-प्रणालीत संकेतबद्ध
  करतो. स्वयंसिद्धके परस्परविरोधी नाहीत, ही गोष्ट आपल्याला लगेच पडताळून
  पाहायची असते. ती परस्परविरोधी असतील, तर त्या नियमांचे पालन करणाऱ्या
  संकल्पना असूच शकत नाहीत आणि आपण विसंगत उपपत्ती मांडण्याचा प्रयत्न केलेला
  असतो. $\Gamma$ चे प्रतिमान शोधून असे घडत नाही हे आपण पडताळू शकतो. आणि
  आपल्या उपपत्तीची प्रतिमाने असतील, तर त्यांचा तपास करण्यासाठी तसेच प्रतिमाने
  रचण्यासाठी आपण तार्किक पद्धती वापरू शकतो.
\item स्वयंसिद्धकांचे स्वातंत्र्य हाही एक महत्त्वाचा प्रश्न आहे. एखादे
  स्वयंसिद्धक प्रत्यक्षात इतरांचा परिणाम असल्यामुळे अनावश्यक असू शकते.
  $\Gamma$ मधील स्वयंसिद्धक~$\varphi$ अनावश्यक आहे हे आपण
  $\Gamma \setminus \{\varphi\} \Entails \varphi$ सिद्ध करून दाखवू शकतो. तसेच,
  $(\Gamma \setminus \{\varphi\}) \cup \{\lnot \varphi\}$ पूर्ततायोग्य आहे हे
  दाखवून ते स्वयंसिद्धक अनावश्यक नाही हेही सिद्ध करू शकतो. उदाहरणार्थ,
  भूमितीचे समांतर गृहीतक इतर स्वयंसिद्धकांपासून स्वतंत्र आहे हे अशाच प्रकारे
  दाखवले गेले.
\item उपपत्तीतील संकल्पनांची परिभाष्यता हाही आणखी एक महत्त्वाचा प्रश्न आहे:
  भाषेच्या निवडीवर उपपत्तीची प्रतिमाने कशाची बनलेली असतात हे ठरते. पण
  उपपत्तीचा प्रत्येक पैलू तिच्या प्रतिमानांत स्वतंत्रपणे दर्शवावाच लागतो असे
  नाही. उदाहरणार्थ, प्रत्येक क्रमसंबंध~$\le$ त्याच्याशी संबंधित काटेकोर
  क्रमसंबंध~$<$ ठरवतो---एक दिला असेल, तर दुसरा परिभाषित करता येतो. म्हणून
  असा क्रम असलेल्या उपपत्तीच्या प्रतिमानात त्याच्याशी संबंधित काटेकोर
  क्रमसंबंध \emph{वेगळ्याने} असणे आवश्यक नाही. सर्वसाधारणपणे, एक संबंध
  इतरांच्या आधारे नेमका कधी परिभाषित करता येतो? एखादा संबंध इतरांच्या
  आधारे परिभाषित करणे कधी अशक्य असते (आणि म्हणून तो भाषेच्या आदिम
  संकल्पनांत जोडावा लागतो)?
\end{enumerate}
\end{explain}












\section{संरचनेचे गुणधर्म व्यक्त करणे}\label{fol:mat:exs:sec}

\begin{explain}
फलने आणि संबंध यांच्यावरील अटी व्यक्त करणे, किंवा अधिक सर्वसाधारणपणे,
रचनेतील फलने आणि संबंध या अटींची पूर्ति करतात हे व्यक्त करणे अनेकदा
उपयुक्त आणि महत्त्वाचे असते. उदाहरणार्थ, एखाद्या भाषेसाठीच्या ज्या
संरचना आपल्याला विधेयांचा जो अर्थ ``पकडायचा'' आहे तो
पकडतात, त्या तसे न करणाऱ्या रचनांपासून वेगळ्या ओळखण्याचे मार्ग आपल्याला
हवे असतात. अर्थात, कोणती संरचना आपल्याला ``अभिप्रेत'' आहेत हे
ठरवण्याचे पूर्ण स्वातंत्र्य आपल्याला आहे. उदाहरणार्थ, आपण असे ठरवू शकतो
की विधेय~$\le$ चे अर्थनिर्धारण हा क्रमसंबंध असलाच पाहिजे; किंवा
आपल्याला फक्त~$\Lang L$ च्या अशा अर्थनिर्धारणांत रस आहे की ज्यांचा प्रांत
संचांचा बनलेला आहे आणि ज्यांत~$\Obj \in$ चे अर्थनिर्धारण ``चा
घटक आहे'' या संबंधाने केलेले आहे. पण भाषेतील वाक्ये वापरून
आपण हे करू शकतो का? दुसऱ्या शब्दांत, संरचना~$\Struct M$ वरील
कोणत्या अटी आपण~$\Struct M$ च्या भाषेतील वाक्य (किंवा कदाचित
वाक्यांचा संच) वापरून व्यक्त करू शकतो? काही अटी आपल्याला व्यक्त करता
येणार नाहीत. उदाहरणार्थ, फक्त $\Domain M = \Nat$ असलेल्या
संरचना~$\Struct M$ मध्येच सत्य असे~$\Lang L_A$ चे कोणतेही वाक्य
नाही. ``प्रांतात फक्त नैसर्गिक संख्या आहेत'' हे आपण व्यक्त करू शकत नाही.
पण संरचनेचे काही ``रचनात्मक गुणधर्म'' आपण कदाचित व्यक्त करू शकतो.
संरचनेचे कोणते गुणधर्म आपण वाक्ये वापरून व्यक्त करू शकतो?
किंवा, दुसऱ्या प्रकारे विचारल्यास, एखादे वाक्य (किंवा
वाक्यांचा संच) सत्य करणाऱ्या रचना म्हणून आपण
संरचनेच्या कोणत्या संग्रहांचे वर्णन करू शकतो?
\end{explain}

\begin{defn}[संचाचे प्रतिमान]
भाषा~$\Lang L$ मधील वाक्यांचा संच~$\Gamma$ असू द्या. प्रत्येक
$\varphi \in \Gamma$ साठी $\Sat{M}{\varphi}$ असेल, तर संरचना~$\Struct M$
\emph{हे~$\Gamma$ चे प्रतिमान आहे} असे म्हणतो.
\end{defn}

\begin{ex}
वाक्य $\lforall[x][x \le x]$ हे~$\Struct M$ मध्ये तेव्हा आणि केवळ तेव्हाच
सत्य असते, जेव्हा $\Assign{\le}{M}$ हा परावर्ती संबंध असतो. वाक्य
$\lforall[x][\lforall[y][((x \le y \land y \le x) \lif x = y)]]$ हे
~$\Struct M$ मध्ये तेव्हा आणि केवळ तेव्हाच सत्य असते, जेव्हा
$\Assign{\le}{M}$ प्रतिसममित असतो. वाक्य
$\lforall[x][\lforall[y][\lforall[z][((x \le y \land y \le z)
      \lif x \le z)]]]$ हे~$\Struct M$ मध्ये तेव्हा आणि केवळ तेव्हाच सत्य
असते, जेव्हा $\Assign{\le}{M}$ संक्रमक असतो. म्हणून पुढील संचाची प्रतिमाने
\begin{align*}
\{\quad &\lforall[x][x \le x], \\
   & \lforall[x][\lforall[y][((x \le y \land y \le
    x) \lif x = y)]], \\
   &\lforall[x][\lforall[y][\lforall[z][((x \le y
      \land y \le z) \lif x \le z)]]] \quad \}
\end{align*}
म्हणजे नेमकी ती रचना होत ज्यांत~$\Assign{\le}{M}$ हा परावर्ती,
प्रतिसममित आणि संक्रमक, म्हणजेच अंशतः क्रम, असतो. त्यामुळे ही वाक्ये
आपण \emph{अंशतः क्रमांच्या प्रथम-क्रम उपपत्तीची} स्वयंसिद्धके म्हणून घेऊ
शकतो.
\end{ex}












\section{प्रथम-क्रम उपपत्तींची उदाहरणे}\label{fol:mat:the:sec}

\begin{ex}
भाषा~$\Lang L_<$ मधील काटेकोर रेषीय क्रमांची उपपत्ती पुढील संचाने
स्वयंसिद्धकांद्वारे निरूपित होते:
\begin{align*}
\{\quad & \lforall[x][\lnot x < x], \\
& \lforall[x][\lforall[y][((x < y \lor y <
    x) \lor x = y)]], \\
& \lforall[x][\lforall[y][\lforall[z][((x < y
      \land y < z) \lif x < z)]]] \quad \}
\end{align*}
ही उपपत्ती अभिप्रेत संरचना पूर्णपणे पकडते: प्रत्येक काटेकोर
रेषीय क्रम या स्वयंसिद्धक-प्रणालीचे प्रतिमान असतो; आणि उलट, $R$ हा
संच~$X$ वरील रेषीय क्रम असेल, तर $\Domain M = X$ आणि
$\Assign{<}{M} = R$ असलेली रचना~$\Struct M$ या उपपत्तीचे प्रतिमान असते.
\end{ex}

\begin{ex}
$\Obj 1$ (स्थिरांक) आणि $\cdot$ (द्विस्थानी फलन) असलेल्या
भाषेतील गटांची उपपत्ती पुढील वाक्यांनी स्वयंसिद्धकांद्वारे निरूपित होते:
\begin{align*}
& \lforall[x][\eq[(x \cdot \Obj 1)][x]]\\
& \lforall[x][\lforall[y][\lforall[z][\eq[(x \cdot (y \cdot z))][((x
          \cdot y) \cdot z)]]]]\\
& \lforall[x][\lexists[y][\eq[(x \cdot y)][\Obj 1]]]
\end{align*}
\end{ex}

\begin{ex}
पेआनो अंकगणिताची उपपत्ती अंकगणिताची भाषा~$\Lang L_A$ मधील पुढील
वाक्यांनी स्वयंसिद्धकांद्वारे निरूपित होते.
\begin{align*}
& \lforall[x][\lforall[y][(\eq[x'][y'] \lif \eq[x][y])]]\\
& \lforall[x][\eq/[\Obj 0][x']]\\
& \lforall[x][\eq[(x + \Obj 0)][x]]\\
& \lforall[x][\lforall[y][\eq[(x + y')][(x + y)']]]\\
& \lforall[x][\eq[(x \times \Obj 0)][\Obj 0]]\\
& \lforall[x][\lforall[y][\eq[(x \times y')][((x \times y) + x)]]]\\
& \lforall[x][\lforall[y][(x < y \liff \lexists[z][\eq[(z' + x)][y])]]]\\
\intertext{आणि पुढील रूपाची सर्व वाक्ये}
& (\varphi(\Obj 0) \land \lforall[x][(\varphi(x) \lif \varphi(x'))]) \lif \lforall[x][\varphi(x)]
\end{align*}
शेवटच्या रूपाची अनंतपणे अनेक वाक्ये असल्यामुळे ही स्वयंसिद्धक-प्रणाली
अनंत आहे. शेवटच्या रूपाला \emph{विगमन-रूपबंध} म्हणतात. (प्रत्यक्षात,
विगमन-रूपबंधाचे स्वरूप आपण येथे दाखवले आहे त्यापेक्षा थोडे अधिक गुंतागुंतीचे
आहे.)

शेवटचे स्वयंसिद्धक हे~$<$ ची \emph{स्पष्ट व्याख्या} आहे.
\end{ex}

\begin{ex}
शुद्ध संचांची उपपत्ती गणिताच्या पायामध्ये (आणि गणिताच्या तत्त्वज्ञानातही)
महत्त्वाची भूमिका बजावते. एखाद्या संचाचे सर्व घटक हेही शुद्ध संच
असतील, तर तो संच शुद्ध असतो. म्हणून रिक्त संच शुद्ध मानला जातो; पण एखाद्या
संचात संच नसलेली कोणतीही वस्तू घटक म्हणून असेल, तर तो संच शुद्ध
नसतो. म्हणजे शुद्ध संच हे फक्त रिक्त संचापासून, कोणतेही ``मूलघटक'' न
वापरता बनवलेले संच होत; मूलघटक म्हणजे स्वतः संच नसलेल्या वस्तू.

पुढील वाक्ये शुद्ध संचांच्या उपपत्तीसाठी स्वयंसिद्धक-प्रणाली मानता येतील:
\begin{align*}
& \lexists[x][\lnot \lexists[y][y \in x]]\\
& \lforall[x][\lforall[y][(\lforall[z](z \in x \liff z \in y) \lif
\eq[x][y])]]\\
& \lforall[x][\lforall[y][\lexists[z][\lforall[u][(u \in z \liff
(\eq[u][x] \lor \eq[u][y]))]]]]\\
& \lforall[x][\lexists[y][\lforall[z][(z \in y \liff \lexists[u][(z \in
        u \land u \in x)])]]]\\
\intertext{आणि पुढील रूपाची सर्व वाक्ये} &
\lexists[x][\lforall[y][(y \in x \liff \varphi(y))]]
\end{align*}
पहिले स्वयंसिद्धक सांगते की कोणतेही घटक नसलेला एक संच आहे (म्हणजे
$\emptyset$ अस्तित्वात आहे); दुसरे सांगते की संच विस्तारात्मक आहेत; तिसरे
सांगते की कोणत्याही संच~$X$ आणि~$Y$ साठी संच~$\{X, Y\}$ अस्तित्वात आहे;
आणि चौथे सांगते की कोणत्याही संच~$X$ साठी संच~$\cup X$ अस्तित्वात आहे,
जिथे~$\cup X$ हा~$X$ च्या सर्व घटकांचा संयोग आहे.

सर्वांत शेवटी उल्लेखलेल्या वाक्यांना एकत्रितपणे
\emph{अनौपचारिक गुणधर्माधारित संचरचना-रूपबंध} म्हणतात. तो मूलतः असे सांगतो
की प्रत्येक~$\varphi(x)$ साठी संच~$\Setabs{x}{\varphi(x)}$ अस्तित्वात आहे---म्हणून
पहिल्या नजरेत ते सत्य, उपयुक्त आणि कदाचित आवश्यकही वाटणारे स्वयंसिद्धक आहे.
त्याला ``अनौपचारिक'' म्हणतात कारण, प्रत्यक्षात, त्यामुळे ही उपपत्ती
अपूर्ततायोग्य होते: $\varphi(y)$ म्हणून $\lnot y \in y$ घेतल्यास पुढील
वाक्य मिळते:
\[\lexists[x][\lforall[y][(y \in x \liff \lnot y \in y)]]
\]
आणि या वाक्याची कोणत्याही संरचनांमध्ये पूर्ति होत नाही.
\end{ex}

\begin{ex}
\emph{अवयवमीमांसेमध्ये} \emph{अवयवत्वाचा} संबंध हा एक मूलभूत संबंध आहे.
संचांच्या उपपत्तींप्रमाणेच अवयवत्वाच्याही अशा उपपत्ती आहेत, ज्या या
संबंधाविषयीच्या विविध (आणि कधीकधी परस्परविरोधी) संकल्पना
स्वयंसिद्धकांद्वारे निरूपित करतात.

अवयवमीमांसेच्या भाषेत एकच द्विस्थानी विधेयचिन्ह~$\Obj P$ असते आणि
$\Atom{\Obj P}{x, y}$ चा ``अर्थ'' $x$ हा~$y$ चा अवयव आहे असा असतो. हे
अर्थनिर्धारण मनात ठेवले असता, या भाषेतील संरचना ला
\emph{अवयवत्वाची रचना} म्हणतात. अर्थात, एकच द्विस्थानी विधेय असलेली प्रत्येक
रचना या नावास खरोखर पात्र असेलच असे नाही. ``अवयवत्व'' पकडण्याची शक्यता
असण्यासाठी $\Assign{\Obj P}{M}$ ने काही अटींची पूर्ति केली पाहिजे; त्या अटी
आपण अवयवत्वाच्या उपपत्तीची स्वयंसिद्धके म्हणून मांडू शकतो. उदाहरणार्थ,
वस्तूंवरील अवयवत्व हा अंशतः क्रम आहे: प्रत्येक वस्तू स्वतःचाच अवयव असते
(जरी तो \emph{अनुचित} अवयव असला तरी); कोणत्याही दोन भिन्न वस्तू एकमेकींचे
अवयव असू शकत नाहीत; आणि एखाद्या वस्तूच्या अवयवाचा अवयव हा स्वतः त्या
वस्तूचा अवयव असतो. लक्षात घ्या की या अर्थाने ``चा अवयव आहे'' हे ``चा
उपसंच आहे'' यासारखे आहे; पण ते ``चा घटक आहे'' यासारखे नाही, कारण घटक-संबंध
परावर्तीही नाही आणि संक्रमकही नाही.
\begin{align*}
& \lforall[x][\Atom{\Obj P}{x,x}] \\
& \lforall[x][\lforall[y][((\Part{x}{y} \land \Part{y}{x})
      \lif \eq[x][y])]] \\
& \lforall[x][\lforall[y][\lforall[z][((\Part{x}{y} \land
        \Part{y}{z}) \lif \Part{x}{z})]]]\\
\intertext{याखेरीज, कोणत्याही दोन वस्तूंचा अवयवी संयोग असतो (ज्याचे त्या दोन
  वस्तू अवयव असतात आणि जो या बाबतीत न्यूनतम असतो).}  &
\lforall[x][\lforall[y][\lexists[z][\lforall[u][(\Part{z}{u} \liff
        (\Part{x}{u} \land \Part{y}{u}))]]]]
\end{align*}
तत्त्वमीमांसकांनी विचारात घेतलेल्या अवयवत्वाच्या मूलभूत तत्त्वांपैकी ही
फक्त काही तत्त्वे आहेत. मात्र, आधी काही परिभाषित संबंध न मांडता पुढील
तत्त्वांची सूत्रे बनवणे किंवा ती लिहून काढणे लवकरच कठीण होते. उदाहरणार्थ,
अवयवमीमांसेत रस असलेले बहुतेक तत्त्वमीमांसक पुढील तत्त्वही वैध मानतात:
जेव्हा वस्तू~$x$ ला उचित अवयव~$y$ असतो, तेव्हा तिला असा अवयव~$z$ देखील
असतो की त्या अवयवाचा~$y$ शी कोणताही सामाईक अवयव नसतो आणि त्यामुळे~$y$ व
$z$ यांचा अवयवी संयोग~$x$ असतो.
\end{ex}












\section{संरचना मधील संबंध व्यक्त करणे}\label{fol:mat:exr:sec}

\begin{explain}
सूत्रांचा एक प्रमुख उपयोग म्हणजे संरचना~$\Struct M$ मधील
गुणधर्म आणि संबंध हे~$\Struct M$ ची भाषा~$\Lang L$ हिच्या आदिम संकल्पनांच्या
साहाय्याने व्यक्त करणे. याचा अर्थ पुढीलप्रमाणे आहे: $\Struct M$ चा
प्रांत हा वस्तूंचा संच असतो. स्थिरांक, फलने आणि
विधेये यांचे~$\Struct M$ मध्ये अनुक्रमे~$\Domain M$ मधील काही वस्तू,
$\Domain M$ वरील फलने आणि~$\Domain M$ वरील संबंध यांनी अर्थनिर्धारण केलेले
असते. उदाहरणार्थ, $\Obj A^2_0$ हे~$\Lang L$ मध्ये असेल, तर~$\Struct M$
त्याला संबंध~$R = \Assign{{\Obj A^2_0}}{M}$ नेमते. मग सूत्र
$\Atom{\Obj A^2_0}{\Obj v_1, \Obj v_2}$ पुढील अर्थाने नेमका तो संबंध
\emph{व्यक्त करते}: चर-मूल्यांकन~$s$ हे~$\Obj v_1$ ला
$a \in \Domain{M}$ आणि~$\Obj v_2$ ला $b \in \Domain M$ शी संगत करत असेल,
तर
\[Rab \text{\quad तेव्हा आणि केवळ तेव्हाच\quad} \Sat{M}{\Atom{\Obj A^2_0}{\Obj v_1, \Obj v_2}}[s].
\]
येथे चर-मूल्यांकनांचा समावेश करणे आवश्यक आहे हे लक्षात घ्या: आपण फक्त
``$Rab$ iff $\Sat{M}{\Atom{\Obj A^2_0}{a, b}}$'' असे म्हणू शकत नाही, कारण
$a$ आणि~$b$ ही आपल्या भाषेतील चिन्हे नाहीत; ती~$\Domain{M}$ ची
घटक आहेत.

आपल्याकडे केवळ अणुसूत्र सूत्रे नसून, तार्किक संयोजके आणि संख्यापक
वापरून ती एकत्र करता येतात; त्यामुळे अधिक गुंतागुंतीची सूत्रे अशी इतर
संबंध परिभाषित करू शकतात जी~$\Struct M$ मध्ये थेट अंतर्भूत नसतात. ते कसे
करायचे आणि विशेषतः संरचनांमध्ये कोणते संबंध परिभाषित करता येतात,
यात आपल्याला रस आहे.
\end{explain}

\begin{defn}
$\varphi(\Obj v_1,\dots, \Obj v_n)$ हे~$\Lang L$ चे असे सूत्र असू द्या
की ज्यात फक्त~$\Obj v_1$,\dots, $\Obj v_n$ हीच चरे मुक्त येतात; आणि
$\Struct M$ ही~$\Lang L$ साठीची संरचना असू द्या. कोणत्याही अशा
चर-मूल्यांकन~$s$ साठी की $s(\Obj v_i) = a_i$ ($i = 1, \dots, n$),
\[Ra_1\dots a_n \text{\quad तेव्हा आणि केवळ तेव्हाच\quad} \Sat{M}{\Atom{\varphi}{\Obj
    v_1,\dots, \Obj v_n}}[s]
\]
असे असेल, तर आणि केवळ तेव्हाच $\varphi(\Obj v_1,\dots, \Obj v_n)$ हे
$R \subseteq \Domain M^n$ हा \emph{संबंध व्यक्त करते}.
\end{defn}

\begin{ex}
अंकगणिताच्या मानक प्रतिमान~$\Struct N$ मध्ये सूत्र $\Obj
v_1 < \Obj v_2 \lor \eq[\Obj v_1][\Obj v_2]$ हे~$\Nat$ वरील~$\le$
संबंध व्यक्त करते. सूत्र $\eq[\Obj v_2][\Obj v_1']$ हे उत्तरवर्ती
संबंध व्यक्त करते; म्हणजे, $m$ हा~$n$ चा उत्तरवर्ती असेल तेव्हा
$Rnm$ असणारा संबंध~$R \subseteq \Nat^2$. सूत्र
$\eq[\Obj v_1][\Obj v_2']$ हे पूर्ववर्ती संबंध व्यक्त करते. सूत्रे
$\lexists[\Obj v_3][(\eq/[\Obj v_3][\Obj 0] \land
  \eq[\Obj v_2][(\Obj v_1 + \Obj v_3)])]$ आणि $\lexists[\Obj
  v_3][\eq[(\Obj v_1 + {\Obj v_3}')][\Obj v_2]]$ ही दोन्ही~$\Obj <$
संबंध व्यक्त करतात.

याचा अर्थ अंकगणिताच्या भाषेत विधेयचिन्ह~$<$ प्रत्यक्षात अनावश्यक आहे; ते
परिभाषित करता येते.
\end{ex}

\begin{explain}
ही कल्पना फक्त विशिष्ट संरचनांमध्येच रंजक नाही; एखादे अभिप्रेत
प्रतिमान किंवा प्रतिमाने वर्णन करण्यासाठी आपण भाषा वापरतो तेव्हा, म्हणजे
उपपत्तींचा विचार करताना, ती सर्वसाधारणपणेही रंजक ठरते. अशा उपपत्तींमध्ये
मूलभूत चिन्हे म्हणून अनेकदा फक्त काही विधेये असतात, पण ती ज्या
प्रांताचे वर्णन करतात त्यात इतर अनेक संबंध महत्त्वाची भूमिका बजावतात. या
इतर संबंधांना भाषेतील मूलभूत विधेयांचे अर्थनिर्धारण करणाऱ्या
संबंधांनी सुव्यवस्थित रीतीने व्यक्त करता येत असेल, तर आपण त्यांना त्या
भाषेत \emph{परिभाषित} करू शकतो असे म्हणतो.
\end{explain}

\begin{prob}
पुढील संबंध परिभाषित करणारी~$\Lang L_A$ मधील सूत्रे शोधा:
\begin{enumerate}
\item $n$ हे~$i$ आणि~$j$ यांच्या दरम्यान आहे;
\item $n$ ने~$m$ ला निःशेष भाग जातो (म्हणजे $m$ ही~$n$ ची पटीत संख्या आहे);
\item $n$ ही अविभाज्य संख्या आहे (म्हणजे $1$ आणि~$n$ यांखेरीज कोणत्याही
  संख्येने~$n$ ला निःशेष भाग जात नाही).
\end{enumerate}
\end{prob}

\begin{prob}
सूत्र $\varphi(\Obj v_1, \Obj v_2)$ हे संरचना~$\Struct M$ मध्ये संबंध
$R \subseteq \Domain M^2$ व्यक्त करते असे समजा. पुढील संबंध व्यक्त करणारी
सूत्रे शोधा:
\begin{enumerate}
\item $R$ चा व्युत्क्रम~$R^{-1}$;
\item सापेक्ष गुणाकार~$R \mid R$;
\end{enumerate}
$R$ चे संक्रमक संवरण~$R^+$ व्यक्त करण्याचा मार्ग तुम्हाला शोधता येतो का?
\end{prob}

\begin{prob}
फक्त एक द्विस्थानी विधेयचिन्ह~$<$ असलेली भाषा~$\Lang{L}$ असू द्या (इतर
कोणतीही स्थिरांक, फलने किंवा विधेये नाहीत---अर्थात
~$\eq$ सोडून). $\Domain{N} = \Nat$ आणि
$\Assign{<}{N} = \Setabs{\tuple{n,m}}{n < m}$ असणारी रचना~$\Struct{N}$
असू द्या. पुढील विधाने सिद्ध करा:
\begin{enumerate}
\item $\{ 0 \}$ हा~$\Struct{N}$ मध्ये परिभाष्य आहे;
\item $\{ 1 \}$ हा~$\Struct{N}$ मध्ये परिभाष्य आहे;
\item $\{ 2 \}$ हा~$\Struct{N}$ मध्ये परिभाष्य आहे;
\item प्रत्येक $n \in \Nat$ साठी संच~$\{ n \}$ हा~$\Struct{N}$ मध्ये
  परिभाष्य आहे;
\item $\Domain{N}$ चा प्रत्येक सांत उपसंच~$\Struct{N}$ मध्ये परिभाष्य आहे;
\item $\Domain{N}$ चा प्रत्येक सांत-पूरक उपसंच~$\Struct{N}$ मध्ये परिभाष्य
  आहे (जिथे $X \subseteq \Nat$ हा सांत-पूरक असतो तेव्हा आणि केवळ तेव्हाच,
  जेव्हा $\Nat \setminus X$ हा सांत असतो).
\end{enumerate}
\end{prob}












\section{संचांची उपपत्ती}\label{fol:mat:set:sec}

जवळजवळ संपूर्ण गणित संचांच्या उपपत्तीमध्ये विकसित करता येते. या
उपपत्तीमध्ये गणित विकसित करताना अनेक गोष्टी कराव्या लागतात. प्रथम,
संबंध~$\in$ साठी स्वयंसिद्धकांचा संच लागतो. वेगवेगळ्या अनेक
स्वयंसिद्धक-प्रणाली विकसित झाल्या असून, त्यांतील~$\in$ चे गुणधर्म कधीकधी
परस्परविरोधी असतात. निवडीच्या स्वयंसिद्धकासह झर्मेलो--फ्रेंकेल संच
उपपत्ती म्हणून ओळखली जाणारी स्वयंसिद्धक-प्रणाली~$\Log{ZFC}$ विशेष उठून
दिसते: ती आतापर्यंत सर्वाधिक वापरली आणि अभ्यासली गेलेली प्रणाली आहे, कारण
गणितज्ञांना ज्या जवळजवळ सर्व गोष्टी सिद्ध करता याव्यात अशी अपेक्षा असते,
त्या सिद्ध करण्यास तिची स्वयंसिद्धके पुरेशी ठरतात. पण हे स्थापित करण्यापूर्वी
गणितज्ञांना व्यक्त करायच्या सर्व गोष्टी आपण \emph{व्यक्त} तरी कशा करू शकतो,
हे प्रथम स्पष्ट करणे आवश्यक आहे. सुरुवातीलाच, भाषेत कोणतीही स्थिरांक
किंवा फलने नसतात; त्यामुळे विशिष्ट संचांविषयी (उदाहरणार्थ,
$\emptyset$ किंवा~$\Nat$), संचांवरील क्रियांविषयी (उदाहरणार्थ,
$X \cup Y$ आणि~$\Pow{X}$), इतकेच नव्हे तर संबंध आणि फलने यांसारख्या
संचांखेरीज इतर गोष्टींचा समावेश असलेल्या रचनांविषयी आपण बोलू शकतो हेही
पहिल्या नजरेत अस्पष्ट वाटते.

सुरुवातीला, ``चा घटक आहे'' हाच काही आपल्याला महत्त्वाचा असलेला एकमेव संबंध
नाही; ``चा उपसंच आहे'' हा जवळजवळ तितकाच महत्त्वाचा वाटतो. पण ``चा घटक
आहे'' याच्या आधारे ``चा उपसंच आहे'' हे आपण \emph{परिभाषित} करू शकतो. त्यासाठी
संच उपपत्तीच्या भाषेत असे सूत्र~$\varphi(x, y)$ शोधावे लागते की संचांची
जोडी~$\tuple{X, Y}$ त्याची पूर्ति तेव्हा आणि केवळ तेव्हाच करते, जेव्हा
$X \subseteq Y$. पण~$X$ हा~$Y$ चा उपसंच तेव्हा आणि केवळ तेव्हाच असतो,
जेव्हा~$X$ चे सर्व घटक हे~$Y$ चेही घटक असतात. म्हणून
पुढील सूत्राने~$\subseteq$ ची व्याख्या करता येते:
\[\lforall[z][(z \in x \lif z \in y)]
\]
आता, एखाद्या सूत्रात~$\subseteq$ संबंध वापरायचा असेल तेव्हा त्याऐवजी हे
सूत्र वापरता येईल ($x$ आणि~$y$ यांच्या जागी योग्य पदे घालून आणि आवश्यक
असल्यास बद्ध चर~$z$ चे नाव बदलून). उदाहरणार्थ, संचांची विस्तारात्मकता
म्हणजे कोणतेही संच~$x$ आणि~$y$ हे परस्परांचे उपसंच असतील, तर~$x$ आणि~$y$
एकच संच असले पाहिजेत. हे
$\lforall[x][\lforall[y][((x \subseteq y \land y
    \subseteq x) \lif x = y)]]$ ने व्यक्त करता येते; किंवा वरच्या व्याख्येने
$\subseteq$ बदलल्यास पुढील सूत्राने:
\[\lforall[x][\lforall[y][((\lforall[z][(z \in x \lif z \in y)] \land
    \lforall[z][(z \in y \lif z \in x)]) \lif x = y)]].
\]
हे खरे तर~$\Log{ZFC}$ च्या स्वयंसिद्धकांपैकी एक, ``विस्तारात्मकतेचे
स्वयंसिद्धक,'' आहे.

$\emptyset$ साठी कोणतेही स्थिरांक नाही, पण ``$x$ रिक्त आहे'' हे
$\lnot \lexists[y][y \in x]$ ने व्यक्त करता येते. मग ``$\emptyset$
अस्तित्वात आहे'' हे वाक्य~$\lexists[x][\lnot \lexists[y][y \in
    x]]$ बनते. हे~$\Log{ZFC}$ चे आणखी एक स्वयंसिद्धक आहे. (विस्तारात्मकतेच्या
स्वयंसिद्धकानुसार फक्त एकच रिक्त संच असतो, हे लक्षात घ्या.) संच उपपत्तीच्या
भाषेत~$\emptyset$ विषयी बोलायचे असेल तेव्हा आपण ते ``रिक्त असलेला एक संच
अस्तित्वात आहे आणि \dots'' असे लिहू. उदाहरणार्थ, $\emptyset$ हा प्रत्येक संचाचा
उपसंच आहे ही गोष्ट पुढीलप्रमाणे व्यक्त करता येईल:
\[\lexists[x][(\lnot\lexists[y][y \in x] \land \lforall[z][x \subseteq
    z])]
\]
येथे अर्थातच $x \subseteq z$ हेसुद्धा त्याच्या व्याख्येने बदलावे लागेल.

$X \cup Y$ आणि~$\Pow{X}$ यांसारख्या संचांवरील क्रियांविषयी बोलण्यासाठी
अशीच युक्ती वापरावी लागते. संच उपपत्तीच्या भाषेत फलनचिन्हे नसतात, पण
$X \cup Y = Z$ आणि~$\Pow{X} = Y$ हे फलनीय संबंध पुढील सूत्रांनी व्यक्त
करता येतात:
\begin{align*}
& \lforall[u][((u \in x \lor u \in y) \liff u \in z)]\\
& \lforall[u][(u \subseteq x \liff u \in y )]
\end{align*}
कारण $X \cup Y$ चे घटक म्हणजे नेमके~$X$ चे घटक किंवा
~$Y$ चे घटक असलेले संच, आणि~$\Pow{X}$ चे घटक म्हणजे नेमके~$X$
चे उपसंच होत. पण त्यामुळे आपण $x \cup y$ किंवा~$\Pow{x}$ ही पदे असल्यासारखी
वापरू शकत नाही; $X \cup Y = Z$ आणि~$\Pow{X} = Y$ हे संबंध परिभाषित करणारी
संपूर्ण सूत्रे एवढीच वापरता येतात. प्रत्यक्षात, या संबंधांची कधी तरी
पूर्ति होते, म्हणजे संयोग आणि घातसंच नेहमी अस्तित्वात असतात, हे आपल्याला
अजून माहीत नाही. उदाहरणार्थ, वाक्य
$\lforall[x][\lexists[y][\Pow{x} = y]]$ हे~$\Log{ZFC}$ चे आणखी एक
स्वयंसिद्धक (घातसंचाचे स्वयंसिद्धक) आहे.

आता क्रमित जोड्या किंवा फलने यांविषयी कसे बोलायचे? येथे क्रमित जोड्या आणि
फलने हे विशिष्ट प्रकारचे संच म्हणून कसे मानता येतात, हे स्पष्ट करावे लागते.
क्रमित जोडी~$\tuple{x, y}$ ची व्याख्या संच~$\{\{x\}, \{x, y\}\}$ म्हणून
करणे हा एक मार्ग आहे. पण आधीप्रमाणेच या संचाला नाव देणारे फलन
आपण समाविष्ट करू शकत नाही; फक्त $\tuple{x, y} = z$, म्हणजे
$\{\{x\}, \{x, y\}\} = z$, हा संबंध परिभाषित करता येतो:
\[\lforall[u][(u \in z \liff (\lforall[v][(v \in u \liff v = x)] \lor
  \lforall[v][(v \in u \liff (v = x \lor v = y))]))]
\]
याचा अर्थ~$z$ चे घटक~$u$ म्हणजे नेमके असे संच की ज्यांचा एकमेव
घटक हा~$x$ असतो किंवा ज्यांचे एकमेव घटक हे~$x$ आणि~$y$
असतात (दुसऱ्या शब्दांत, जे संच~$\{x\}$ शी किंवा~$\{x, y\}$ शी एकरूप
असतात). हे मिळाल्यानंतर आणखी गोष्टीही सांगता येतात; उदाहरणार्थ,
$X \times Y = Z$:
\[\lforall[z][(z \in Z \liff \lexists[x][\lexists[y][(x \in
        X \land y \in Y \land \tuple{x, y} = z)]])]
\]

फलन~$f \colon X \to Y$ याचा संबंध~$f(x) = y$, म्हणजे जोड्यांचा संच
$\Setabs{\tuple{x,y}}{f(x) = y}$, म्हणून विचार करता येतो. मग संच~$f$ हे
$X$ पासून~$Y$ पर्यंतचे फलन आहे असे आपण तेव्हा म्हणू शकतो, जेव्हा (a) तो
$\subseteq X \times Y$ असा संबंध आहे, (b) तो सर्वत्र परिभाषित आहे,
म्हणजे प्रत्येक $x \in X$ साठी असा काही $y \in Y$ आहे की
$\tuple{x, y} \in f$, आणि (c) तो एकमूल्यी आहे, म्हणजे
$\tuple{x, y}, \tuple{x, y'} \in f$ असताना $y = y'$ (कारण फलनांची मूल्ये
एकमेव असली पाहिजेत). म्हणून ``$f$ हे~$X$ पासून~$Y$ पर्यंतचे फलन आहे'' हे
पुढीलप्रमाणे लिहिता येते:
\begin{align*}
 \lforall[u][(u \in f \lif {}] & \lexists[x][\lexists[y][(x \in X \land y \in
      Y \land \tuple{x, y} = u)]]) \land {}\\
 \lforall[x][(x \in X \lif {}] &
      (\lexists[y][(y \in Y \land \mathrm{maps}(f, x, y))] \land {}\\
& (\lforall[y][\lforall[y'][((\mathrm{maps}(f, x, y) \land
    \mathrm{maps}(f, x, y')) \lif y = y')]]))
\end{align*}
येथे $\mathrm{maps}(f, x, y)$ हे $\lexists[v][(v \in f \land
  \tuple{x, y} = v)]$ चे संक्षिप्त रूप आहे (हे सूत्र ``$f(x) = y$''
व्यक्त करते).

आता, उदाहरणार्थ, $f\colon X \to Y$ हे एकास-एक आहे हे व्यक्त करणेही
कठीण नाही:
\begin{multline*}
 f \colon X \to Y \land \lforall[x][\lforall[x'][((x \in X \land x' \in
     X \land {}]] \\
 \lexists[y][(\mathrm{maps}(f, x, y) \land \mathrm{maps}(f,
        x', y))]) \lif x = x')
\end{multline*}
फलन~$f\colon X \to Y$ हे एकास-एक तेव्हा आणि केवळ तेव्हाच असते, जेव्हा
$f$ हे $x, x' \in X$ या दोन्हींना एकाच~$y$ शी संगत करत असल्यास $x = x'$.
हे सूत्र~$\mathrm{inj}(f, X, Y)$ असे संक्षिप्तपणे लिहिल्यास, आता आपण संच
उपपत्तीच्या भाषेत कँटरच्या प्रमेयासारखी साधी नसलेली गोष्टही मांडू शकतो:
$\Pow{X}$ पासून~$X$ मध्ये जाणारे कोणतेही एकास-एक फलन नाही:
\[\lforall[X][\lforall[Y][(\Pow{X} = Y \lif
    \lnot\lexists[f][\mathrm{inj}(f, Y, X)])]]
\]

प्रत्येक परिभाषक गुणधर्मासाठी संचाचे अस्तित्व हमीने देणारे आणखी एक
स्वयंसिद्धक संच उपपत्तीला आवश्यक आहे असे वाटू शकते. $\varphi(x)$ हे मुक्त
चर~$x$ असलेले संच उपपत्तीचे सूत्र असेल, तर पुढील वाक्य विचारात घेता
येते:
\[\lexists[y][\lforall[x][(x \in y \liff \varphi(x))]].
\]
हे वाक्य सांगते की असा संच~$y$ आहे ज्याचे घटक म्हणजे नेमके
ते सर्व~$x$ होत जे~$\varphi(x)$ ची पूर्ति करतात. या रूपबंधाला ``गुणधर्माधारित
संचरचनेचे तत्त्व'' म्हणतात. ते फार उपयुक्त वाटते; पण दुर्दैवाने ते विसंगत
आहे. $\varphi(x) \ident \lnot x \in x$ घ्या; मग गुणधर्माधारित संचरचनेचे तत्त्व
पुढील विधान सांगते:
\[\lexists[y][\lforall[x][(x \in y \liff x \notin x)]],
\]
म्हणजे स्वतःचे घटक नसलेल्या सर्व संचांच्या संचाचे अस्तित्व ते
सांगते. असा संच अस्तित्वात असू शकत नाही---ही रसेलची विरोधापत्ती आहे.
प्रत्यक्षात, $\Log{ZFC}$ मध्ये या तत्त्वाची मर्यादित---आणि सुसंगत---आवृत्ती,
म्हणजे विभक्तीकरण तत्त्व, समाविष्ट आहे:
\[\lforall[z][\lexists[y][\lforall[x][(x \in y \liff (x \in z \land
      \varphi(x))]]].
\]

\begin{prob}
पुढील निष्पन्नता दाखवणारी निष्पत्ती देऊन गुणधर्माधारित संचरचनेचे
तत्त्व विसंगत आहे, हे दाखवा:
\[\lexists[y][\lforall[x][(x \in y \liff x \notin x)]] \Proves \lfalse.
\]
प्रथम $(A \lif \lnot A) \land (\lnot A \lif A)
\Proves \lfalse$ दाखवल्यास मदत होईल.
\end{prob}











\section{संरचनेचे आकारमान व्यक्त करणे}\label{fol:mat:siz:sec}

\begin{explain}
भाषेतील अतार्किक चिन्हे न वापरताही रचनांचे काही गुणधर्म व्यक्त करता येतात.
उदाहरणार्थ, अशी वाक्ये आहेत की ती संरचनांमध्ये तेव्हा आणि केवळ
तेव्हाच सत्य असतात, जेव्हा त्या संरचनेच्या प्रांत मध्ये
किमान, जास्तीत जास्त किंवा नेमके दिलेल्या संख्येइतके~$n$ घटक असतात.
\end{explain}

\begin{prop}
वाक्य
\begin{multline*}
  \varphi_{\ge n} \ident \lexists[x_1][\lexists[x_2][\dots\lexists[x_n][{}]]]\\
  \begin{aligned}
  (\eq/[x_1][x_2] \land {}
  \eq/[x_1][x_3] \land \eq/[x_1][x_4] \land \dots \land \eq/[x_1][x_n] \land {}\\
\eq/[x_2][x_3] \land \eq/[x_2][x_4] \land \dots \land {} \eq/[x_2][x_n] \land {} \\
\vdots\\
\eq/[x_{n-1}][x_n])
\end{aligned}
\end{multline*}
हे संरचना~$\Struct M$ मध्ये तेव्हा आणि केवळ तेव्हाच सत्य असते,
जेव्हा $\Domain M$ मध्ये किमान $n$ घटक असतात. परिणामी,
$\Sat{M}{\lnot \varphi_{\ge n+1}}$ तेव्हा आणि केवळ तेव्हाच लागू होते, जेव्हा
$\Domain M$ मध्ये जास्तीत जास्त~$n$ घटक असतात.

\end{prop}

\begin{prop}
वाक्य
\begin{multline*}
\varphi_{= n} \ident \lexists[x_1][\lexists[x_2][\dots\lexists[x_n][{}]]] \\
\begin{aligned}
  (\eq/[x_1][x_2] \land {}
  \eq/[x_1][x_3] \land \eq/[x_1][x_4] \land \dots \land \eq/[x_1][x_n] \land {}\\
\eq/[x_2][x_3] \land \eq/[x_2][x_4] \land \dots \land {} \eq/[x_2][x_n] \land {} \\
\vdots\\
\eq/[x_{n-1}][x_n] \land {} \\
\lforall[y][(\eq[y][x_1] \lor \dots \lor \eq[y][x_n]]))
\end{aligned}
\end{multline*}
हे संरचना~$\Struct M$ मध्ये तेव्हा आणि केवळ तेव्हाच सत्य असते,
जेव्हा $\Domain M$ मध्ये नेमके $n$ घटक असतात.
\end{prop}

\begin{prop}
एखादी संरचना अनंत तेव्हा आणि केवळ तेव्हाच असते, जेव्हा ती पुढील
संचाचे प्रतिमान असते:
\[\{\varphi_{\ge 1}, \varphi_{\ge 2}, \varphi_{\ge 3}, \dots \}.
\]
\end{prop}

असे एकही पूर्णतः तार्किक वाक्य नाही की जे~$\Struct M$ मध्ये तेव्हा आणि केवळ
तेव्हाच सत्य असेल, जेव्हा $\Domain M$ अनंत असेल. पण अशी अतार्किक
विधेये असलेली वाक्ये देता येतात ज्यांची फक्त अनंत प्रतिमाने
असतात (जरी प्रत्येक अनंत संरचना त्यांचे प्रतिमान नसते). सांत रचना
असण्याचा गुणधर्म आणि अगणनीय रचना असण्याचा गुणधर्म
वाक्यांच्या अनंत संचानेही व्यक्त करता येत नाही. ही तथ्ये संहतता आणि
लोव्हेनहाइम--स्कोलेम प्रमेयांवरून निष्पन्न होतात.



